\originalchapter{官方讲义对应与范围}
官方四份讲义共60个物理页, 含每份的 MIT OCW 来源页及章标题页, 数学正文页数不同于课程讲次. 下列对应基于实际读到的讲义, 不由课程名称推测. 第27--29讲的逆函数和水平集内容在原件有重复, 本书合并到曲面定义章.
\begin{longtable}{p{.18\textwidth}p{.39\textwidth}p{.34\textwidth}}
\toprule
官方讲义与讲次 & 原件实际内容 & 本书处理\\
\midrule\endhead
第1章, 1--3 & 平面曲率、弧长、Frenet、曲线基本定理、密切圆、水平集曲率 & 第1章补条件与证明, 配椭圆和隐式算例.\\
第1章, 4 & Minkowski 平面曲线 & 不纳入欧氏曲线主线.\\
第1章, 5--6 & 高维 Frenet 标架及三维曲率、挠率 & 第3章作三维特化, 补曲线基本定理与螺线判据.\\
第1章, 7--10 & 圆映射次数、绕数、旋转数、Jordan、转角、凸性、Whitney、四顶点 & 第2章取切角、转角、严格正曲率凸性和四顶点; 不给 Whitney 分类或一般正则同伦理论.\\
第2章, 11--15 & 正定内积、超曲面基本形式、形算子、主曲率、旋转及切线曲面 & 第4--6章作二维曲面特化, 明示平均曲率因子差异.\\
第2章, 16--18,21 & Christoffel、Gauss 方程、曲率内蕴性、移动标架、特殊坐标 & 第7章逐式证明, 增补 Codazzi 和曲面基本定理局部存在性. 正常坐标在第9章构造.\\
第2章, 19--20,22--23 & 高维外代数曲率、刚性、Minkowski、余维推广、抽象度量 & 高维与 Lorentz 推广不纳入; 第12章使用二维圆盘内蕴度量.\\
第3章, 24--30 & 水平集、切空间、定向、图像化、紧致超曲面、Hadamard & 第4--5、11章二维特化; 逆隐函数为先修, 凸性结论给几何证明并明示覆盖工具.\\
第3章, 31--35 & 面积、映射次数、Gauss 映射总曲率、带奇点标架、Gauss--Bonnet、组合曲率 & 第4、11章取曲面面积和 Gauss Jacobian. 用局部公式和三角剖分证明整体 Gauss--Bonnet. 不系统讲一般次数、Hopf 高维版本或多面体重复内容.\\
第4章, 36--39 & 测地线、旋转曲面、变分、Hamilton、内蕴距离、双曲距离、Schwarz--Pick & 第8--9章证明测地线与最短性. 第12章证明圆盘距离. 不扩展 Hamilton 动力系统或复分析 Schwarz--Pick.\\
第4章, 40--41 & Busemann/CAT、测地曲率、带边界 Gauss--Bonnet、测地三角形 & 第10--11章完整几何公式及角点; 度量曲率空间不纳入.\\
\bottomrule
\end{longtable}
本册增加的教学材料包括图像曲率完整计算、悬链面和螺旋面、法向面积变分、Clairaut 关系细节、Gauss 引理、闭嵌入曲面的最短线存在证明以及36道练习的完整解答. 这些内容是对选定主线的重构与推导, 不伪称每一题均为原课程作业或原作者的逐字陈述. 本册不讨论一般流形的系统拓扑、向量丛特征类、高维 Riemann 几何、完整负曲率嵌入障碍、Jacobi 场与完整二阶变分理论; 这一区域的缺失不以“流形入门”替代.

\par\Needspace{8\baselineskip}
\originalsection{官方作业的覆盖与缺口}
10份公开Homework每份2个物理页, 共20页, 含10个题面页. 逐题实际计数为3、2、3、5、3、4、3、2、4、3, 共32个顶层题. 原件四题带有10个正式标号子问, 属于上述32题内部. 作业1--10按原次序附在基础正文与36题解答之后, 并保留原分值、条件和指涉. 30个顶层题的全部可恢复要求得到独立解答, 其中原题错误的两个命题给出反例与修正版. 作业3.3的折线定义、总曲率与Hopf部分已解, 但其最后所引Proposition 6.3没有可恢复陈述; 作业9.4仅引用未定位的Lemma 28.3, 因而保留缺口, 不猜题.

作业6.1的商业教材方法未复制, 改为明确署作独立推导的自足翻转隆起构造. 作业7.2的基本形式用词及原点归一化条件、作业10.1的讲次编号、作业10.2的有符号角动量不等式均保留原文并批注. 新增作业还补齐所用外幂、球面次数积分、奇点标架及受约束粒子方程; 这不等于系统展开整个高维课程. 公开目录未提供官方答案或考试原卷, 本册全部新增解答为AI独立编写.
